\documentclass[11pt]{article}
\usepackage[margin=1in]{geometry}
\usepackage{mathpazo}
\usepackage{amsmath,amssymb}
\usepackage{microtype}
\usepackage{booktabs}
\usepackage[dvipsnames]{xcolor}
\usepackage[most]{tcolorbox}
\usepackage{enumitem}
\usepackage{hyperref}
\hypersetup{colorlinks=true,linkcolor=MidnightBlue,urlcolor=MidnightBlue}
\usepackage{cleveref}
\urlstyle{tt}
\def\UrlBreaks{\do\_\do/\do.}
\newcommand{\file}[1]{\nolinkurl{#1}}
\newcommand{\code}[1]{\texttt{#1}}
\newcommand{\mhh}{m_{hh}}

\definecolor{trainable}{HTML}{1B6CA8}
\definecolor{physics}{HTML}{2E933C}
\definecolor{warn}{HTML}{C0392B}

\newtcolorbox{keyidea}[1][]{enhanced,breakable,colback=physics!6,colframe=physics!70!black,
  fonttitle=\bfseries,title={#1},boxrule=0.6pt,arc=2pt,left=6pt,right=6pt}
\newtcolorbox{aside}[1][]{enhanced,breakable,colback=trainable!5,colframe=trainable!70!black,
  fonttitle=\bfseries,title={#1},boxrule=0.6pt,arc=2pt,left=6pt,right=6pt}
\newtcolorbox{watch}[1][]{enhanced,breakable,colback=warn!5,colframe=warn!70!black,
  fonttitle=\bfseries,title={#1},boxrule=0.6pt,arc=2pt,left=6pt,right=6pt}

\setlist{itemsep=2pt,topsep=3pt}

\title{\textbf{Three checks of the \texttt{dihiggs} data and test pipeline}\\[4pt]
\large Plan, for approval before the checks are run}
\date{October 28, 2026}

\begin{document}
\maketitle

\begin{abstract}
\noindent
This note sets out how I propose to check that (1) every HEFT parameter point
in the POWHEG file is unique, (2) the 1\,000 holdout spectra written by
\code{heft\_prepare\_data.ipynb} are disjoint from the data used by
\code{heft\_train.ipynb}, and (3) \code{heft\_test.ipynb} correctly compares the
POWHEG histograms with the predictions of the trained model. For each check I
state the question, why it matters, the method, and what counts as a pass. I also
list things I noticed while reading the code; these are \emph{observations to be
tested}, not results. No checks have been run yet.
\end{abstract}

% ===========================================================================
\section{Ground rules}

\begin{itemize}
  \item \textbf{What is checked.} The repository as of commit \code{2cc3bc5}
  (``more testing''): the three notebooks, the three files in \file{data/}, and
  the trained model in \file{code/runs/L1Loss/}.
  \item \textbf{Read-only.} No file in the repository is modified. In
  particular, \code{heft\_prepare\_data.ipynb} is \emph{not} rerun, because it
  reshuffles the spectra without a fixed random seed. Rerunning it would
  produce new, different data files and destroy the very thing being checked.
  \item \textbf{Where.} Checks 1 and 2 need only pandas and numpy and can run on
  your computer. Check 3 needs PyTorch, so it runs in my workspace on copies of the
  files.
  \item \textbf{Deliverables.} The checks become a notebook,
  \file{code/heft_checks.ipynb}, which prints PASS or FAIL for each test, so
  they can be rerun after every change and later become a step in the workflow.
  The results go in a short PDF report.
\end{itemize}

\begin{aside}[A note on comparing numbers]
\code{heft\_prepare\_data.ipynb} converts the Wilson coefficients and cross
sections to 32-bit floats (\code{float32}) before writing them. The test file stores
these as short decimals, while the training file stores the same values after
conversion to 64-bit floats, written with more digits. The same number can
therefore \emph{look} different in the two files. Wherever values from different
files are compared, I convert both to \code{float32} first. Equal \code{float32}
values then compare exactly, with no tolerance needed.
\end{aside}

% ===========================================================================
\section{Check 1: every parameter point in the POWHEG file is unique}

\paragraph{Question.} Does each row of
\file{data/powheg_total_gauss_params_13TeV.csv.gz} correspond to a different point
$\theta = (\kappa_\lambda, c_t, c_{tt}, c_{ggh}, c_{gghh})$?

\paragraph{Why it matters.} A duplicated point could end up in both the training
and the holdout sets. The holdout would then partly test memory rather than
generalisation, even with perfect bookkeeping. A duplicated row is also a sign that
the file was assembled incorrectly.

\paragraph{Method.}
\begin{enumerate}[label=\textbf{1\alph*.}]
  \item \textbf{Exact duplicates.} Read the file (4\,999 rows expected) and count
  rows whose five coefficients are identical. Do this both as stored and after
  conversion to \code{float32}, which is the precision the pipeline actually uses.
  Two distinct points could in principle become equal on conversion.
  \item \textbf{Near duplicates.} Exact equality can miss points that differ only
  by rounding. Divide each coefficient by its standard deviation, so that the five
  axes are comparable, and compute the distance between every pair of points:
  $4\,999 \times 4\,998/2 \approx 12.5$ million pairs, a few seconds with numpy. I will report
  the smallest distances and list any pair closer than $10^{-6}$, roughly the
  precision of the stored numbers.
  \item \textbf{Duplicate spectra.} Check whether any two rows have identical
  cross sections in all 100 bins. If two different $\theta$ shared a spectrum, that
  would point to a copying error in producing the file.
  \item \textbf{Sanity check of the closest pairs.} For the few closest pairs of
  points, compare their spectra. Nearby points should have similar spectra. A large
  difference would mean that either the parameters or the spectra are attached to
  the wrong row.
\end{enumerate}

\paragraph{Pass.} No exact duplicates, either as stored or in \code{float32}; no
pair closer than $10^{-6}$; no duplicate spectra. The distribution of
nearest-neighbour distances is reported for information.

% ===========================================================================
\section{Check 2: the 1\,000 holdout spectra are disjoint from the training data}

\subsection{How the notebook builds the two sets}

Reading \code{heft\_prepare\_data.ipynb}:
\begin{enumerate}
  \item Each POWHEG row gets an \code{ID} equal to its original row number,
  0--4\,998. The rows are then shuffled (\code{df.sample(frac=1)}).
  \item The first 1\,000 shuffled rows are written to
  \file{heft_wide_gauss_testdata_13TeV.csv} (the holdout), keeping their \code{ID}.
  \item The remaining 3\,999 rows are diced into one row per bin, with their
  \code{ID}. Zero cross sections are dropped, the rows are shuffled, and they are
  written to \file{heft_wide_gauss_traindata_13TeV.csv.gz}.
  \item \code{heft\_train.ipynb} reads only the training file and splits it into
  training (85\%), validation (2\%) and test (13\%) parts.
\end{enumerate}
By construction, the two sets are disjoint \emph{if both files were written by the
same run} of the notebook.

\begin{watch}[What could go wrong]
\begin{itemize}
  \item \textbf{Files from different runs.} The shuffle has no fixed seed, so every
  run makes a different split. If the training file came from one run and the
  holdout file from another, about a quarter of the holdout spectra would be in the
  training data. The git history has separate commits ``update traindata'' and
  ``update testdata'', so this needs checking rather than assuming.
  \item \textbf{A model trained on older data.} Before the change from 500 to 1\,000
  holdout spectra, the training file contained 4\,499 spectra. A model trained on
  that file would have seen up to 500 of today's holdout spectra. The model in
  \file{runs/L1Loss} must have been trained on the \emph{current} training file.
\end{itemize}
\end{watch}

\subsection{Method}
\begin{enumerate}[label=\textbf{2\alph*.}]
  \item \textbf{Counts.} The holdout file has 1\,000 rows with 1\,000 distinct
  \code{ID}s. The training file's \code{ID}s form a set of 3\,999. The total number
  of training rows equals $3\,999 \times 80$ minus the number of zero bins in those
  spectra. The notebook comment records 318\,115.
  \item \textbf{Disjoint by ID.} The two sets of \code{ID}s have no element in
  common, and together they make up exactly $\{0, 1, \dots, 4\,998\}$.
  \item \textbf{Disjoint by value.} \code{ID}s could be mislabelled, so repeat the
  test using the parameter values themselves. No $\theta$ (in \code{float32}) in
  the training file may appear in the holdout file. Together with Check 1, this is
  independent of the bookkeeping.
  \item \textbf{Each file is a faithful copy of the POWHEG file.}
  \begin{itemize}
    \item Every holdout row must equal POWHEG row \code{ID}: the same five
    coefficients and the same 80 cross sections, bins 17--96.
    \item For every training row, the coefficients must equal those of POWHEG row
    \code{ID}. The mass value $x$ must be a bin centre, $x = (k - 17 + \tfrac12)/100$
    for an integer $k$ between 17 and 96, and $\sigma$ must equal POWHEG row
    \code{ID}, bin $k$.
    \item Each \code{ID} must contribute exactly its number of non-zero bins.
  \end{itemize}
  This catches misaligned columns or a wrong bin-to-mass mapping, as well as mixed
  files.
  \item \textbf{No leak inside the training notebook.} Confirm, by reading the
  code, that \code{heft\_train.ipynb} opens no file other than the training file.
  Also confirm that the \code{drwho} \code{Dataset} and \code{DataLoader} classes
  only slice the arrays they are given and never read from disk.
  \item \textbf{The trained model used the current training file.}
  \begin{itemize}
    \item The sizes recorded in \file{runs/L1Loss/dihiggs_config.yaml} must
    equal those computed in the notebook from the current file:
    \code{int(0.85 N)}, \code{int(0.13 N)} and the remainder, with
    $N = 318\,115$.
    \item The parameter file must be newer than the commit that last changed the
    training file.
  \end{itemize}
\end{enumerate}

\paragraph{Pass.} All of 2a--2f hold. If 2f fails, the bookkeeping may be correct
but the current model cannot be fairly tested on the current holdout, and would
have to be retrained first.

\begin{aside}[{Something worth knowing, though not a failure}]
The 13\% ``test'' split inside \code{heft\_train.ipynb} consists of \emph{bins},
not spectra. Its rows come from the same 3\,999 parameter points as the training
rows, just from other bins. It therefore tests interpolation in $\mhh$, not
generalisation to new $\theta$. Only the 1\,000 holdout spectra test the
property needed for reweighting. I will state this in the report and not count it
as a failure.
\end{aside}

% ===========================================================================
\section{Check 3: the test notebook compares like with like}

\paragraph{Question.} For each holdout spectrum, does \code{heft\_test.ipynb}
overlay the right POWHEG histogram and the model's prediction for the \emph{same}
$\theta$, at the right masses, in the same units?

\subsection{Method}
\begin{enumerate}[label=\textbf{3\alph*.}]
  \item \textbf{Code review.} Read \code{get\_spectra} and \code{plot\_spectra}
  line by line and confirm four things:
  \begin{itemize}
    \item the bin columns selected are 17--96, in order;
    \item the mass values are the same bin centres the training file used;
    \item the input columns are in the order the model expects:
    \code{mhh, klambda, CT, CTT, CGGH, CGGHH};
    \item the histogram bins line up with the centres.
  \end{itemize}
  \item \textbf{Independent recomputation.} Run the notebook's own
  \code{get\_spectra} on all 1\,000 holdout spectra. Compare the result with a
  second, independent calculation:
  \begin{itemize}
    \item the POWHEG spectrum is taken directly from the POWHEG file by \code{ID},
    not from the holdout file;
    \item the prediction comes from the exported numpy copy of the model, which
    takes $\mhh$ in GeV and builds its inputs with different code.
  \end{itemize}
  The POWHEG spectra must agree exactly (in \code{float32}). The predictions must
  agree to about $10^{-6}$ relative, the difference between the notebook's 32-bit
  and the copy's 64-bit arithmetic.
  \item \textbf{Pairing survives sorting.} The notebook sorts the spectra by total
  cross section before plotting. For every sorted entry, check that its
  coefficients, POWHEG spectrum and prediction still belong together, by looking up
  the coefficients in the POWHEG file.
  \item \textbf{The histogram reproduces the data.} Confirm that
  \code{ax.hist(x, bins=80, range=(0, 0.8), weights=f)} puts each bin centre in its
  own bin, so that the drawn histogram heights are exactly the POWHEG values.
  \item \textbf{A spot check by eye.} Redraw one page of 20 panels independently
  and compare it with the notebook's figure for the same page.
\end{enumerate}

\paragraph{Pass.} 3a--3e hold. Separately, for information, I will summarise how
well the model agrees with POWHEG over the 1\,000 spectra, using the notebook's
measure and a per-bin relative one (below).

\begin{watch}[{Observations from reading the code, to be confirmed}]
\begin{itemize}
  \item \textbf{What ``RMS$/\sigma$'' measures.} The panels show
  $\sqrt{\text{bias}^2 + \text{variance}}\,/\,\sigma_\text{total}$, computed over the 80
  bins. That is the root-mean-square difference per \emph{bin} divided by the
  \emph{total} cross section, which is the sum of 80 bins. It is therefore
  about 80 times smaller than the error relative to a typical bin. This is not a bug,
  but the label may mislead, and a per-bin relative measure is easier to interpret.
  \item \textbf{The coefficient plot covers only part of the range.}
  \code{plot\_coeffs} draws $a_i$ for $0 \le x \le 0.5$, that is
  $255 \le \mhh \le 1005$~GeV, while the model is trained up to 1455~GeV.
  \item \textbf{Minor fragilities.} \code{get\_spectra} uses the global
  \code{spectra} rather than its argument \code{df}. This is harmless here,
  because they are the same object. Sorting a list of tuples that contain numpy
  arrays would fail if two spectra had exactly the same total, which is unlikely.
\end{itemize}
\end{watch}

% ===========================================================================
\section{Summary of tests}

\begin{center}
\small
\begin{tabular}{lll}
\toprule
test & what is compared & pass if\\
\midrule
1a & the five coefficients, row against row & no exact duplicates\\
1b & standardised distances between all pairs & none below $10^{-6}$\\
1c & 100-bin spectra, row against row & no identical spectra\\
1d & spectra of the closest pairs & similar (by eye)\\
\midrule
2a & row and \code{ID} counts & 1\,000 holdout; 3\,999 training IDs\\
2b & holdout \code{ID}s against training \code{ID}s & disjoint; union is $0$--$4\,998$\\
2c & holdout $\theta$ against training $\theta$ (\code{float32}) & disjoint\\
2d & both files against the POWHEG file & identical values\\
2e & files read by \code{heft\_train.ipynb} & training file only\\
2f & model configuration and dates against the training file & sizes match; model newer\\
\midrule
3a & notebook code & correct bins, masses, input order\\
3b & notebook against an independent calculation & exact (POWHEG), $\sim10^{-6}$ (model)\\
3c & entries after sorting & all still paired correctly\\
3d & drawn histogram against POWHEG values & identical\\
3e & one page redrawn independently & same picture\\
\bottomrule
\end{tabular}
\end{center}

\begin{keyidea}[Decisions for you]
\begin{enumerate}
  \item Is a near-duplicate threshold of $10^{-6}$ (in standardised units) what
  you want, or do you have a physically motivated scale in mind?
  \item May I commit \file{code/heft_checks.ipynb} to the repository, so the
  checks become part of the workflow?
  \item If 2f shows the model predates the current training file, should I stop
  and report, or retrain and then run Check 3?
\end{enumerate}
\end{keyidea}

\end{document}
